\section*{Solid-State Drive (SSD) Technology}
A Solid-State Drive (SSD) is a non-volatile storage device that stores digital data using NAND flash memory
rather than magnetic platters. Unlike Hard Disk Drives (HDDs), SSDs contain no moving mechanical parts,
enabling significantly faster data access, lower latency, reduced power consumption, and improved durability.
SSDs have become the preferred storage solution for modern laptops, desktops, servers, and embedded systems
because they provide high performance while maintaining reliability and energy efficiency [1, 3].
SSDs communicate with the host computer through storage interfaces such as Serial ATA (SATA) and PCI
Express (PCIe). SATA SSDs use the AHCI protocol originally designed for HDDs, whereas newer SSDs may
use the NVMe protocol over PCIe for even higher performance. Regardless of interface, the fundamental
storage mechanism of SSDs remains based on NAND flash memory [4].

\subsubsection{Physical Infrastructure}


\begin{enumerate}[label=\roman*.]
    \item NAND Flash Memory
    The primary storage medium in an SSD consists of NAND flash memory chips. These chips permanently store
    data without requiring continuous electrical power [3].
   \item  SSD Controller
    The SSD controller acts as the brain of the drive. It manages read and write operations, error correction, garbage
    collection, wear leveling, logical-to-physical address mapping, and communication with the operating system
    [4, 6].
    \item DRAM Cache (Optional)
    Many SSDs include a DRAM cache that stores frequently accessed mapping information and temporary data,
    reducing access latency and improving performance [4].
    \item Interface Connector
    Consumer SSDs commonly connect through SATA connectors or M.2 slots, depending on the drive type.
    Enterprise SSDs may also use U.2 or PCIe expansion card interfaces [4].
\end{enumerate}